\section{Introduction}
\label{sec:b-introduction}

An autonomous experimentation agent can fail in ways that a warning score
foresees but no single response repairs.  The same warning can justify
continuing, invoking an independent verifier, or rolling back to a known-good
state; the options differ in benefit, in resource cost, and in whether they
damage a checkpoint that would have succeeded on its own
\citep{interventionparadox2026,harnessr12026}.  Branching from a saved
checkpoint is the natural way to observe those differences: restore the state,
execute each option, compare \citep{calibrationnotcontrol2026,causalagentreplay2026,replaygap2026}.
But a branch contrast is a measurement before it is a decision, and the
measurement has preconditions --- the option must be applicable, and every
branch must begin from an equivalent executable state --- that are the subject
of separate work \citep{zhang2026ca}.  This paper takes those preconditions as
an explicit contract and asks the decision question: \emph{how should a
controller allocate feasible interventions under resource and harm
constraints?}

\paragraph{What is already established, and what this paper adds.}
The closest prior work, \citet{calibrationnotcontrol2026}, names
\emph{intervention advantage} as the decision quantity, formalises the
prediction/control mismatch as target error, and introduces \emph{prefix
branching}: candidate actions are executed from an identical trajectory state
and the resulting action-conditioned outcomes support an intervention decision.
That paper already establishes that prediction and intervention value come
apart, and the repository record of its inspected scope shows a fixed action
menu, reconstruction checks with mismatch rejection, and a repeated
shallow-prefix probe.  \citet{interventionparadox2026} show that a critic with
strong offline discrimination can still degrade agent performance, decomposing
the effect into recovery and disruption.  \citet{harnessr12026} learn
lifecycle-wide executable interventions and report that many valid patches
reduce batch success.  The ordinal condition under which an outcome ranking
coincides with an effect ranking is established in the budgeted-triage and
uplift literatures \citep{okati2021,fernandezloria2022,kamran2024}; we cite it
rather than re-claim it.

What those results do not supply is a decision procedure for the situation in
which the three constraints interact: the action may be \emph{unavailable} at
the offered checkpoint, verification has a \emph{cost} that competes with other
checkpoints, and an intervention can \emph{harm} a run that was on track.  A
controller that ranks by failure risk or by a utility deficit ignores
feasibility and cost; a controller that intervenes everywhere ignores harm; and
a controller that ranks by estimated advantage still needs a rule for how much
of its budget to spend and what evidence licenses the spend.  The contribution
of this paper is to pose that allocation problem precisely, specify the branch
data that make it answerable, and define a comparison that can falsify the
proposed increment.

\paragraph{The proposed contribution, in falsifiable form.}
At an online-reachable decision checkpoint, with one continue/verify decision
under a frozen continuation, we compare two allocation policies at matched
information, capacity and branch budget:
\begin{enumerate}[label=(B\arabic*),leftmargin=2.6em]
\item an \emph{outcome-based} policy that spends its budget on the highest
  predicted failure risk (and, when utility is primary, the largest predicted
  utility deficit); and
\item an \emph{action-conditioned} policy that spends its budget on the largest
  estimated verification advantage, with feasibility, cost and harm handled
  explicitly.
\end{enumerate}
The study is designed to falsify the claim that the action-conditioned
allocation is better under the studied law: if the paired held-out gain
difference $D_O(q)$ of \cref{eq:b-primary-estimand} is indistinguishable from
zero at the primary capacity, the increment does not survive its own test.  The
conservative-selection component is falsified if it removes the benefit it is
meant to protect, or if its bound cannot be defended at the stated scope.  Both
outcomes are reportable; neither is a failure of the design.

\paragraph{What this paper is not.}
It is not a claim that prediction and intervention value come apart --- that is
\citet{interventionparadox2026} and \citet{calibrationnotcontrol2026}.  It is
not a new general identification theorem for replay; the measurement conditions
are handled by a separate contract and cited.  It is not sequential
reinforcement learning: the controller acts once, at one checkpoint, and the
continuation is frozen.  Recovery is not in the first study; it appears only as
an optional third action with a defined precondition, rollback boundary and
continuation contract, and sequential decision-making is future work
(\cref{sec:b-bridge}).  And it is not an empirical result: the paper reports a
design, a contract and an engineering smoke path, with every performance number
marked \emph{planned} (\cref{sec:b-experiments}).

\paragraph{Organisation.}
\Cref{sec:b-problem} fixes the one-decision problem and the online checkpoint
rule.  \Cref{sec:b-regimes} defines the continue and verify regimes.
\Cref{sec:b-data} gives the dataset-generation contract and the runtime
qualification boundary.  \Cref{sec:policy} presents the action-value baseline
and the proposed conservative allocation procedure.  \Cref{sec:b-theory}
records the retained decision theory and its attribution.
\Cref{sec:b-experiments} is the experimental plan with hypotheses, baselines
and figure slots.  \Cref{sec:b-statistics} fixes outcomes, splits, estimands
and inference.  \Cref{sec:b-limitations} states limitations, and
\cref{sec:related} positions the work against the nearest priors.
\Cref{sec:b-bridge} sketches the bridge to recovery and sequential control.
